\documentclass[12pt]{article}
\usepackage[utf8]{inputenc}
\usepackage[T1]{fontenc}
\usepackage{hyperref}
\title{Basque Language Resources for K-12 Curriculum Development}
\author{Compiled by Claude Sonnet 5.5 (claude-sonnet-5-5) for Andr\'as Kornai}
\date{}
\begin{document}
\maketitle

\section{Introduction}
Basque (Euskara; ISO 639-3 eus; Glottocode basq1248) is a language isolate, official in the Basque Autonomous Community and Navarre (Spain), and used as a language of instruction. The search results confirm a mature immersion school system, the ikastolas, but returned no official curriculum documents or textbooks. Search was limited to a few queries, so this list is far from complete; the Basque Government education department, Ikastolen Elkartea (Spain) and Seaska (France) are obvious places to look next, though their sites were not in the results.

\section{Primary Educational Resources}
\begin{itemize}
\item Ikastolen Elkartea document (PDF, in the search results; contents not verified beyond its appearing on ikastola.eus): \url{https://ikastola.eus/sites/default/files/article/41781/file/BIHAR\%20OSOA\%20FRANTSESEZ.pdf}
\item Article analysing French-teaching materials used in the ikastolas, University of the Basque Country: \url{https://ekoizpen-zientifikoa.ehu.eus/documentos/60a856d105fac20a35c959b4}
\item Euskal Kultura, page on Boiseko Ikastola (Basque preschool in Idaho): \url{https://euskalkultura.eus/english/resources/boiseko-ikastola}
\end{itemize}

\section{Linguistic and Reference Materials}
\begin{itemize}
\item Basque Government, Euskara Batua (standard Basque) page: \url{https://www.euskadi.eus/eusko-jaurlaritza//contenidos/informacion/euskara_batua/eu_7419/euskara_batua.html}
\item Euskaltzaindia research document (PDF; contents not verified): \url{https://www.euskaltzaindia.eus/dok/ikerbilduma/26112.pdf}
\item Elhuyar dictionary, press release on 30 years of the dictionary (Basque--Spanish, Basque--French, Basque--English): \url{https://www.elhuyar.eus/en/press-room/dictionary-elhuyar-30-years-and-75-million-consultations}
\item Elhuyar dictionary article: \url{https://zientzia.eus/artikuluak/dictionnaire-elhuyar-hiztegia/en/}
\item Argia report on the Euskaltzaindia grammar committee dispute: \url{https://argia.eus/en/argia-astekaria/2395/euskaltzaindiko-gramatika-batzordearen-eta-ibon-sarasolaren-dimisioak}
\end{itemize}

\section{Religious and Community Resources}
None found specifically; the ikastola community is the main community resource (see below).

\section{Digital and Audio Resources}
The Elhuyar articles above describe the online dictionary with example sentences and audio (user recordings from Forvo, and text-to-speech); I did not see the dictionary site itself in the results.

\section{Curriculum Development Context}
\begin{itemize}
\item Overview of the ikastola school system and its history (Buber's Basque Page): \url{https://buber.net/Basque/2019/12/08/basque-fact-of-the-week-ikastolas-the-basque-schools/}
\item Argia on French-state recognition of Basque-only immersion at a school: \url{https://argia.eus/en/albistea/frantziako-gobernuak-euskara-hutsezko-irakaskuntza-onartu-du-hiriburukjo-baste-quieta-ikastetxean}
\item Idaho News, on America's only Basque immersion school: \url{https://idahonews.com/news/local/inside-americas-only-basque-immersion-school}
\item Boise Dev, 2026 article on Boiseko Ikastola: \url{https://boisedev.com/news/2026/06/11/for-nearly-3-decades-boise-preschool-remains-only-one-of-its-kind-in-the-us/}
\end{itemize}
Basque is already taught and used as a medium from preschool to the baccalauréat in Basque-medium schools, so the work here is selecting and cataloguing existing materials rather than creating a curriculum from nothing.
\end{document}